\frfilename{mt1a.tex}
\versiondate{3.1.04}
\copyrightdate{1996}

\def\chaptername{Lampiran}
\def\sectionname{Pendahuluan}

\gdef\topparagraph{}
\gdef\bottomparagraph{Lampiran untuk Jilid\ 1 {\it pendahuluan}}

\centerline{\bf Lampiran untuk Jilid 1}

\medskip

\centerline{\bf Fakta-Fakta Berguna}

\medskip

Setiap jilid risalah ini akan dilengkapi sebuah lampiran yang memuat uraian
sangat ringkas mengenai bahan yang mungkin sudah pernah dijumpai oleh banyak
pembaca, tetapi mungkin belum oleh sebagian pembaca lain, dan yang relevan
dengan suatu topik yang dibahas dalam jilid tersebut. Untuk jilid pertama ini,
lampirannya singkat, sebagian karena jilidnya sendiri singkat, tetapi terutama
karena pengetahuan analisis yang diperlukan bersifat begitu mendasar sehingga
harus dipelajari dengan semestinya dari buku teks biasa atau dalam mata kuliah
biasa. Meskipun demikian, saya tetap menguraikan beberapa perincian yang
mungkin tidak dibahas dalam sebagian mata kuliah pengantar analisis, yakni
beberapa notasi yang tidak sepenuhnya baku dan teori elementer mengenai
himpunan terhitung (\S1A1), himpunan terbuka dan tertutup dalam ruang Euklides
(\S1A2), serta limit atas dan limit bawah barisan dan fungsi (\S1A3).

\discrpage
